\haddockmoduleheading{Data.Char}
\label{module:Data.Char}
\haddockbeginheader
{\haddockverb\begin{verbatim}
module Data.Char (
    Char, String, isControl, isSpace, isLower, isUpper, isAlpha, isAlphaNum,
    isPrint, isDigit, isOctDigit, isHexDigit, isLetter, isMark, isNumber,
    isPunctuation, isSymbol, isSeparator, isAscii, isLatin1, isAsciiUpper,
    isAsciiLower,
    GeneralCategory(UppercaseLetter,
                    LowercaseLetter,
                    TitlecaseLetter,
                    ModifierLetter,
                    OtherLetter,
                    NonSpacingMark,
                    SpacingCombiningMark,
                    EnclosingMark,
                    DecimalNumber,
                    LetterNumber,
                    OtherNumber,
                    ConnectorPunctuation,
                    DashPunctuation,
                    OpenPunctuation,
                    ClosePunctuation,
                    InitialQuote,
                    FinalQuote,
                    OtherPunctuation,
                    MathSymbol,
                    CurrencySymbol,
                    ModifierSymbol,
                    OtherSymbol,
                    Space,
                    LineSeparator,
                    ParagraphSeparator,
                    Control,
                    Format,
                    Surrogate,
                    PrivateUse,
                    NotAssigned),
    generalCategory, toUpper, toLower, toTitle, digitToInt, intToDigit, ord,
    chr, showLitChar, lexLitChar, readLitChar
  ) where\end{verbatim}}
\haddockendheader

\section{Characters and strings}
\begin{haddockdesc}
\item[\begin{tabular}{@{}l}
data Char
\end{tabular}]
{\haddockbegindoc
The character type \haddockid{Char} represents Unicode codespace
and its elements are code points as in definitions
\href{https://www.unicode.org/versions/Unicode15.0.0/ch03.pdf#G2212}{D9 and D10 of the Unicode Standard}.\par
Character literals in Haskell are single-quoted: \haddocktt{'Q'}, \haddocktt{'Я'} or \haddocktt{'Ω'}.
To represent a single quote itself use \haddocktt{'{\char '134}''}, and to represent a backslash
use \haddocktt{'{\char '134}{\char '134}'}. The full grammar can be found in the section 2.6 of the
\href{https://www.haskell.org/definition/haskell2010.pdf#section.2.6}{Haskell 2010 Language Report}.\par
To specify a character by its code point one can use decimal, hexadecimal
or octal notation: \haddocktt{'{\char '134}65'}, \haddocktt{'{\char '134}x41'} and \haddocktt{'{\char '134}o101'} are all alternative forms
of \haddocktt{'A'}. The largest code point is \haddocktt{'{\char '134}x10ffff'}.\par
There is a special escape syntax for ASCII control characters:\par
{TODO: Table}
\href{https://hackage.haskell.org/package/base/docs/Data-Char.html}{Data.Char} provides utilities to work with \haddockid{Char}.\par}
\end{haddockdesc}
\begin{haddockdesc}
\item[\begin{tabular}{@{}l}
instance Storable Char\\instance Ix Char\\instance Read Char\\instance Show Char\\instance Eq Char\\instance Ord Char\\\end{tabular}]
\end{haddockdesc}
\begin{haddockdesc}
\item[\begin{tabular}{@{}l}
type String = {\char 91}Char{\char 93}
\end{tabular}]
{\haddockbegindoc
\haddockid{String} is an alias for a list of characters.\par
String constants in Haskell are values of type \haddockid{String}.
 That means if you write a string literal like \haddocktt{"hello world"},
 it will have the type \haddocktt{{\char 91}Char{\char 93}}, which is the same as \haddocktt{String}.\par
\textbf{Note:} You can ask the compiler to automatically infer different types
 with the \haddocktt{-XOverloadedStrings} language extension, for example
  \haddocktt{"hello world" :: Text}. See \haddockid{IsString} for more information.\par
Because \haddocktt{String} is just a list of characters, you can use normal list functions
 to do basic string manipulation. See \haddocktt{Data.List} for operations on lists.\par
\subsubsection*{\textbf{Performance considerations}}
\haddocktt{{\char 91}Char{\char 93}} is a relatively memory-inefficient type.
 It is a linked list of boxed word-size characters, internally it looks something like:\par
\begin{quote}
{\haddockverb\begin{verbatim}
╭─────┬───┬──╮  ╭─────┬───┬──╮  ╭─────┬───┬──╮  ╭────╮
│ (:) │   │ ─┼─>│ (:) │   │ ─┼─>│ (:) │   │ ─┼─>│ [] │
╰─────┴─┼─┴──╯  ╰─────┴─┼─┴──╯  ╰─────┴─┼─┴──╯  ╰────╯
        v               v               v
       'a'             'b'             'c'\end{verbatim}}
\end{quote}
The \haddocktt{String} "abc" will use \haddocktt{5*3+1 = 16} (in general \haddocktt{5n+1})
 words of space in memory.\par
Furthermore, operations like \haddockid{(++)} (string concatenation) are \haddocktt{O(n)}
 (in the left argument).\par
For historical reasons, the \haddocktt{base} library uses \haddocktt{String} in a lot of places
 for the conceptual simplicity, but library code dealing with user-data
 should use the \href{https://hackage.haskell.org/package/text}{text}
 package for Unicode text, or the the
 \href{https://hackage.haskell.org/package/bytestring}{bytestring} package
 for binary data.\par}
\end{haddockdesc}
\section{Character classification}
Unicode characters are divided into letters, numbers, marks,
 punctuation, symbols, separators (including spaces) and others (including
 control characters).\par
\begin{haddockdesc}
\item[\begin{tabular}{@{}l}
isControl :: Char -> Bool
\end{tabular}]
{\haddockbegindoc
Selects control characters, which are the non-printing characters of
 the Latin-1 subset of Unicode.\par}
\end{haddockdesc}
\begin{haddockdesc}
\item[\begin{tabular}{@{}l}
isSpace :: Char -> Bool
\end{tabular}]
{\haddockbegindoc
Returns \haddockid{True} for any Unicode space character, and the control
 characters \haddocktt{{\char '134}t}, \haddocktt{{\char '134}n}, \haddocktt{{\char '134}r}, \haddocktt{{\char '134}f}, \haddocktt{{\char '134}v}.\par}
\end{haddockdesc}
\begin{haddockdesc}
\item[\begin{tabular}{@{}l}
isLower :: Char -> Bool
\end{tabular}]
{\haddockbegindoc
Selects lower-case alphabetic Unicode characters (letters).\par
\textbf{Note:} this predicate does \emph{not} work for letter-like characters such as:
 \haddocktt{'ⓐ'} (\haddocktt{U+24D0} circled Latin small letter a) and
 \haddocktt{'ⅳ'} (\haddocktt{U+2173} small Roman numeral four). This is due to selecting only
 characters with the \haddockid{GeneralCategory} \haddockid{LowercaseLetter}.\par
See \haddockid{isLowerCase} for a more intuitive predicate.\par}
\end{haddockdesc}
\begin{haddockdesc}
\item[\begin{tabular}{@{}l}
isUpper :: Char -> Bool
\end{tabular}]
{\haddockbegindoc
Selects upper-case or title-case alphabetic Unicode characters (letters).
 Title case is used by a small number of letter ligatures like the
 single-character form of \emph{Lj}.\par
\textbf{Note:} this predicate does \emph{not} work for letter-like characters such as:
 \haddocktt{'Ⓐ'} (\haddocktt{U+24B6} circled Latin capital letter A) and
 \haddocktt{'Ⅳ'} (\haddocktt{U+2163} Roman numeral four). This is due to selecting only
 characters with the \haddockid{GeneralCategory} \haddockid{UppercaseLetter} or \haddockid{TitlecaseLetter}.\par
See \haddockid{isUpperCase} for a more intuitive predicate. Note that
 unlike \haddockid{isUpperCase}, \haddockid{isUpper} does select \emph{title-case} characters such as
 \haddocktt{'Dž'} (\haddocktt{U+01C5} Latin capital letter d with small letter z with caron) or
 \haddocktt{'ᾯ'} (\haddocktt{U+1FAF} Greek capital letter omega with dasia and perispomeni and
 prosgegrammeni).\par}
\end{haddockdesc}
\begin{haddockdesc}
\item[\begin{tabular}{@{}l}
isAlpha :: Char -> Bool
\end{tabular}]
{\haddockbegindoc
Selects alphabetic Unicode characters (lower-case, upper-case and
 title-case letters, plus letters of caseless scripts and modifiers letters).
 This function is equivalent to \haddockid{isLetter}.\par
This function returns \haddockid{True} if its argument has one of the
 following \haddockid{GeneralCategory}s, or \haddockid{False} otherwise:\par
\vbox{\begin{itemize}
\item
\haddockid{UppercaseLetter}\par
\item
\haddockid{LowercaseLetter}\par
\item
\haddockid{TitlecaseLetter}\par
\item
\haddockid{ModifierLetter}\par
\item
\haddockid{OtherLetter}\par
\end{itemize}}
These classes are defined in the
 \href{http://www.unicode.org/reports/tr44/tr44-14.html#GC_Values_Table}{Unicode Character Database},
 part of the Unicode standard. The same document defines what is
 and is not a "Letter".\par}
\end{haddockdesc}
\begin{haddockdesc}
\item[\begin{tabular}{@{}l}
isAlphaNum :: Char -> Bool
\end{tabular}]
{\haddockbegindoc
Selects alphabetic or numeric Unicode characters.\par
Note that numeric digits outside the ASCII range, as well as numeric
 characters which aren't digits, are selected by this function but not by
 \haddockid{isDigit}. Such characters may be part of identifiers but are not used by
 the printer and reader to represent numbers, e.g., Roman numerals like \haddocktt{\haddocktt{V}},
 full-width digits like \haddocktt{'１'} (aka \haddocktt{'65297'}).\par
This function returns \haddockid{True} if its argument has one of the
 following \haddockid{GeneralCategory}s, or \haddockid{False} otherwise:\par
\vbox{\begin{itemize}
\item
\haddockid{UppercaseLetter}\par
\item
\haddockid{LowercaseLetter}\par
\item
\haddockid{TitlecaseLetter}\par
\item
\haddockid{ModifierLetter}\par
\item
\haddockid{OtherLetter}\par
\item
\haddockid{DecimalNumber}\par
\item
\haddockid{LetterNumber}\par
\item
\haddockid{OtherNumber}\par
\end{itemize}}}
\end{haddockdesc}
\begin{haddockdesc}
\item[\begin{tabular}{@{}l}
isPrint :: Char -> Bool
\end{tabular}]
{\haddockbegindoc
Selects printable Unicode characters
 (letters, numbers, marks, punctuation, symbols and spaces).\par
This function returns \haddockid{False} if its argument has one of the
 following \haddockid{GeneralCategory}s, or \haddockid{True} otherwise:\par
\vbox{\begin{itemize}
\item
\haddockid{LineSeparator}\par
\item
\haddockid{ParagraphSeparator}\par
\item
\haddockid{Control}\par
\item
\haddockid{Format}\par
\item
\haddockid{Surrogate}\par
\item
\haddockid{PrivateUse}\par
\item
\haddockid{NotAssigned}\par
\end{itemize}}}
\end{haddockdesc}
\begin{haddockdesc}
\item[\begin{tabular}{@{}l}
isDigit :: Char -> Bool
\end{tabular}]
{\haddockbegindoc
Selects ASCII digits, i.e. \haddocktt{'0'}..\haddocktt{'9'}.\par}
\end{haddockdesc}
\begin{haddockdesc}
\item[\begin{tabular}{@{}l}
isOctDigit :: Char -> Bool
\end{tabular}]
{\haddockbegindoc
Selects ASCII octal digits, i.e. \haddocktt{'0'}..\haddocktt{'7'}.\par}
\end{haddockdesc}
\begin{haddockdesc}
\item[\begin{tabular}{@{}l}
isHexDigit :: Char -> Bool
\end{tabular}]
{\haddockbegindoc
Selects ASCII hexadecimal digits,
 i.e. \haddocktt{'0'}..\haddocktt{'9'}, \haddocktt{'a'}..\haddocktt{'f'}, \haddocktt{'A'}..\haddocktt{'F'}.\par}
\end{haddockdesc}
\begin{haddockdesc}
\item[\begin{tabular}{@{}l}
isLetter :: Char -> Bool
\end{tabular}]
{\haddockbegindoc
Selects alphabetic Unicode characters (lower-case, upper-case and
 title-case letters, plus letters of caseless scripts and
 modifiers letters). This function is equivalent to
 \haddockid{isAlpha}.\par
This function returns \haddockid{True} if its argument has one of the
 following \haddockid{GeneralCategory}s, or \haddockid{False} otherwise:\par
\vbox{\begin{itemize}
\item
\haddockid{UppercaseLetter}\par
\item
\haddockid{LowercaseLetter}\par
\item
\haddockid{TitlecaseLetter}\par
\item
\haddockid{ModifierLetter}\par
\item
\haddockid{OtherLetter}\par
\end{itemize}}
These classes are defined in the
 \href{http://www.unicode.org/reports/tr44/tr44-14.html#GC_Values_Table}{Unicode Character Database},
 part of the Unicode standard. The same document defines what is
 and is not a "Letter".\par
\subsubsection*{\textbf{Examples}}
Basic usage:\par
\begin{quote}
{\haddockverb\begin{verbatim}
>>> isLetter 'a'
True

>>> isLetter 'A'
True

>>> isLetter 'λ'
True

>>> isLetter '0'
False

>>> isLetter '%'
False

>>> isLetter '♥'
False

>>> isLetter '\31'
False

\end{verbatim}}
\end{quote}
Ensure that \haddockid{isLetter} and \haddockid{isAlpha} are equivalent.\par
\begin{quote}
{\haddockverb\begin{verbatim}
>>> let chars = [(chr 0)..]

>>> let letters = map isLetter chars

>>> let alphas = map isAlpha chars

>>> letters == alphas
True

\end{verbatim}}
\end{quote}}
\end{haddockdesc}
\begin{haddockdesc}
\item[\begin{tabular}{@{}l}
isMark :: Char -> Bool
\end{tabular}]
{\haddockbegindoc
Selects Unicode mark characters, for example accents and the
 like, which combine with preceding characters.\par
This function returns \haddockid{True} if its argument has one of the
 following \haddockid{GeneralCategory}s, or \haddockid{False} otherwise:\par
\vbox{\begin{itemize}
\item
\haddockid{NonSpacingMark}\par
\item
\haddockid{SpacingCombiningMark}\par
\item
\haddockid{EnclosingMark}\par
\end{itemize}}
These classes are defined in the
 \href{http://www.unicode.org/reports/tr44/tr44-14.html#GC_Values_Table}{Unicode Character Database},
 part of the Unicode standard. The same document defines what is
 and is not a "Mark".\par
\subsubsection*{\textbf{Examples}}
Basic usage:\par
\begin{quote}
{\haddockverb\begin{verbatim}
>>> isMark 'a'
False

>>> isMark '0'
False

\end{verbatim}}
\end{quote}
Combining marks such as accent characters usually need to follow
 another character before they become printable:\par
\begin{quote}
{\haddockverb\begin{verbatim}
>>> map isMark "ò"
[False,True]

\end{verbatim}}
\end{quote}
Puns are not necessarily supported:\par
\begin{quote}
{\haddockverb\begin{verbatim}
>>> isMark '✓'
False

\end{verbatim}}
\end{quote}}
\end{haddockdesc}
\begin{haddockdesc}
\item[\begin{tabular}{@{}l}
isNumber :: Char -> Bool
\end{tabular}]
{\haddockbegindoc
Selects Unicode numeric characters, including digits from various
 scripts, Roman numerals, et cetera.\par
This function returns \haddockid{True} if its argument has one of the
 following \haddockid{GeneralCategory}s, or \haddockid{False} otherwise:\par
\vbox{\begin{itemize}
\item
\haddockid{DecimalNumber}\par
\item
\haddockid{LetterNumber}\par
\item
\haddockid{OtherNumber}\par
\end{itemize}}
These classes are defined in the
 \href{http://www.unicode.org/reports/tr44/tr44-14.html#GC_Values_Table}{Unicode Character Database},
 part of the Unicode standard. The same document defines what is
 and is not a "Number".\par
\subsubsection*{\textbf{Examples}}
Basic usage:\par
\begin{quote}
{\haddockverb\begin{verbatim}
>>> isNumber 'a'
False

>>> isNumber '%'
False

>>> isNumber '3'
True

\end{verbatim}}
\end{quote}
ASCII \haddocktt{'0'} through \haddocktt{'9'} are all numbers:\par
\begin{quote}
{\haddockverb\begin{verbatim}
>>> and $ map isNumber ['0'..'9']
True

\end{verbatim}}
\end{quote}
Unicode Roman numerals are "numbers" as well:\par
\begin{quote}
{\haddockverb\begin{verbatim}
>>> isNumber 'Ⅸ'
True

\end{verbatim}}
\end{quote}}
\end{haddockdesc}
\begin{haddockdesc}
\item[\begin{tabular}{@{}l}
isPunctuation :: Char -> Bool
\end{tabular}]
{\haddockbegindoc
Selects Unicode punctuation characters, including various kinds
 of connectors, brackets and quotes.\par
This function returns \haddockid{True} if its argument has one of the
 following \haddockid{GeneralCategory}s, or \haddockid{False} otherwise:\par
\vbox{\begin{itemize}
\item
\haddockid{ConnectorPunctuation}\par
\item
\haddockid{DashPunctuation}\par
\item
\haddockid{OpenPunctuation}\par
\item
\haddockid{ClosePunctuation}\par
\item
\haddockid{InitialQuote}\par
\item
\haddockid{FinalQuote}\par
\item
\haddockid{OtherPunctuation}\par
\end{itemize}}
These classes are defined in the
 \href{http://www.unicode.org/reports/tr44/tr44-14.html#GC_Values_Table}{Unicode Character Database},
 part of the Unicode standard. The same document defines what is
 and is not a "Punctuation".\par
\subsubsection*{\textbf{Examples}}
Basic usage:\par
\begin{quote}
{\haddockverb\begin{verbatim}
>>> isPunctuation 'a'
False

>>> isPunctuation '7'
False

>>> isPunctuation '♥'
False

>>> isPunctuation '"'
True

>>> isPunctuation '?'
True

>>> isPunctuation '—'
True

\end{verbatim}}
\end{quote}}
\end{haddockdesc}
\begin{haddockdesc}
\item[\begin{tabular}{@{}l}
isSymbol :: Char -> Bool
\end{tabular}]
{\haddockbegindoc
Selects Unicode symbol characters, including mathematical and
 currency symbols.\par
This function returns \haddockid{True} if its argument has one of the
 following \haddockid{GeneralCategory}s, or \haddockid{False} otherwise:\par
\vbox{\begin{itemize}
\item
\haddockid{MathSymbol}\par
\item
\haddockid{CurrencySymbol}\par
\item
\haddockid{ModifierSymbol}\par
\item
\haddockid{OtherSymbol}\par
\end{itemize}}
These classes are defined in the
 \href{http://www.unicode.org/reports/tr44/tr44-14.html#GC_Values_Table}{Unicode Character Database},
 part of the Unicode standard. The same document defines what is
 and is not a "Symbol".\par
\subsubsection*{\textbf{Examples}}
Basic usage:\par
\begin{quote}
{\haddockverb\begin{verbatim}
>>> isSymbol 'a'
False

>>> isSymbol '6'
False

>>> isSymbol '='
True

\end{verbatim}}
\end{quote}
The definition of "math symbol" may be a little
 counter-intuitive depending on one's background:\par
\begin{quote}
{\haddockverb\begin{verbatim}
>>> isSymbol '+'
True

>>> isSymbol '-'
False

\end{verbatim}}
\end{quote}}
\end{haddockdesc}
\begin{haddockdesc}
\item[\begin{tabular}{@{}l}
isSeparator :: Char -> Bool
\end{tabular}]
{\haddockbegindoc
Selects Unicode space and separator characters.\par
This function returns \haddockid{True} if its argument has one of the
 following \haddockid{GeneralCategory}s, or \haddockid{False} otherwise:\par
\vbox{\begin{itemize}
\item
\haddockid{Space}\par
\item
\haddockid{LineSeparator}\par
\item
\haddockid{ParagraphSeparator}\par
\end{itemize}}
These classes are defined in the
 \href{http://www.unicode.org/reports/tr44/tr44-14.html#GC_Values_Table}{Unicode Character Database},
 part of the Unicode standard. The same document defines what is
 and is not a "Separator".\par
\subsubsection*{\textbf{Examples}}
Basic usage:\par
\begin{quote}
{\haddockverb\begin{verbatim}
>>> isSeparator 'a'
False

>>> isSeparator '6'
False

>>> isSeparator ' '
True

\end{verbatim}}
\end{quote}
Warning: newlines and tab characters are not considered
 separators.\par
\begin{quote}
{\haddockverb\begin{verbatim}
>>> isSeparator '\n'
False

>>> isSeparator '\t'
False

\end{verbatim}}
\end{quote}
But some more exotic characters are (like HTML's \haddocktt{{\char '46}nbsp;}):\par
\begin{quote}
{\haddockverb\begin{verbatim}
>>> isSeparator '\160'
True

\end{verbatim}}
\end{quote}}
\end{haddockdesc}
\section{Subranges}
\begin{haddockdesc}
\item[\begin{tabular}{@{}l}
isAscii :: Char -> Bool
\end{tabular}]
{\haddockbegindoc
Selects the first 128 characters of the Unicode character set,
 corresponding to the ASCII character set.\par}
\end{haddockdesc}
\begin{haddockdesc}
\item[\begin{tabular}{@{}l}
isLatin1 :: Char -> Bool
\end{tabular}]
{\haddockbegindoc
Selects the first 256 characters of the Unicode character set,
 corresponding to the ISO 8859-1 (Latin-1) character set.\par}
\end{haddockdesc}
\begin{haddockdesc}
\item[\begin{tabular}{@{}l}
isAsciiUpper :: Char -> Bool
\end{tabular}]
{\haddockbegindoc
Selects ASCII upper-case letters,
 i.e. characters satisfying both \haddockid{isAscii} and \haddockid{isUpper}.\par}
\end{haddockdesc}
\begin{haddockdesc}
\item[\begin{tabular}{@{}l}
isAsciiLower :: Char -> Bool
\end{tabular}]
{\haddockbegindoc
Selects ASCII lower-case letters,
 i.e. characters satisfying both \haddockid{isAscii} and \haddockid{isLower}.\par}
\end{haddockdesc}
\section{Unicode general categories}
\begin{haddockdesc}
\item[\begin{tabular}{@{}l}
data GeneralCategory
\end{tabular}]
{\haddockbegindoc
Unicode General Categories (column 2 of the UnicodeData table) in
 the order they are listed in the Unicode standard (the Unicode
 Character Database, in particular).\par
\subsubsection*{\textbf{Examples}}
Basic usage:\par
\begin{quote}
{\haddockverb\begin{verbatim}
>>> :t OtherLetter
OtherLetter :: GeneralCategory

\end{verbatim}}
\end{quote}
\haddockid{Eq} instance:\par
\begin{quote}
{\haddockverb\begin{verbatim}
>>> UppercaseLetter == UppercaseLetter
True

>>> UppercaseLetter == LowercaseLetter
False

\end{verbatim}}
\end{quote}
\haddockid{Ord} instance:\par
\begin{quote}
{\haddockverb\begin{verbatim}
>>> NonSpacingMark <= MathSymbol
True

\end{verbatim}}
\end{quote}
\haddockid{Enum} instance:\par
\begin{quote}
{\haddockverb\begin{verbatim}
>>> enumFromTo ModifierLetter SpacingCombiningMark
[ModifierLetter,OtherLetter,NonSpacingMark,SpacingCombiningMark]

\end{verbatim}}
\end{quote}
\haddockid{Read} instance:\par
\begin{quote}
{\haddockverb\begin{verbatim}
>>> read "DashPunctuation" :: GeneralCategory
DashPunctuation

>>> read "17" :: GeneralCategory
*** Exception: Prelude.read: no parse

\end{verbatim}}
\end{quote}
\haddockid{Show} instance:\par
\begin{quote}
{\haddockverb\begin{verbatim}
>>> show EnclosingMark
"EnclosingMark"

\end{verbatim}}
\end{quote}
\haddockid{Bounded} instance:\par
\begin{quote}
{\haddockverb\begin{verbatim}
>>> minBound :: GeneralCategory
UppercaseLetter

>>> maxBound :: GeneralCategory
NotAssigned

\end{verbatim}}
\end{quote}
\haddockid{Ix} instance:\par
\begin{quote}
{\haddockverb\begin{verbatim}
>>> import GHC.Internal.Data.Ix ( index )

>>> index (OtherLetter,Control) FinalQuote
12

>>> index (OtherLetter,Control) Format
*** Exception: Error in array index

\end{verbatim}}
\end{quote}
\enspace \emph{Constructors}\par
\haddockbeginconstrs
\haddockdecltt{=} & \haddockdecltt{UppercaseLetter} & Lu: Letter, Uppercase \\
\haddockdecltt{|} & \haddockdecltt{LowercaseLetter} & Ll: Letter, Lowercase \\
\haddockdecltt{|} & \haddockdecltt{TitlecaseLetter} & Lt: Letter, Titlecase \\
\haddockdecltt{|} & \haddockdecltt{ModifierLetter} & Lm: Letter, Modifier \\
\haddockdecltt{|} & \haddockdecltt{OtherLetter} & Lo: Letter, Other \\
\haddockdecltt{|} & \haddockdecltt{NonSpacingMark} & Mn: Mark, Non-Spacing \\
\haddockdecltt{|} & \haddockdecltt{SpacingCombiningMark} & Mc: Mark, Spacing Combining \\
\haddockdecltt{|} & \haddockdecltt{EnclosingMark} & Me: Mark, Enclosing \\
\haddockdecltt{|} & \haddockdecltt{DecimalNumber} & Nd: Number, Decimal \\
\haddockdecltt{|} & \haddockdecltt{LetterNumber} & Nl: Number, Letter \\
\haddockdecltt{|} & \haddockdecltt{OtherNumber} & No: Number, Other \\
\haddockdecltt{|} & \haddockdecltt{ConnectorPunctuation} & Pc: Punctuation, Connector \\
\haddockdecltt{|} & \haddockdecltt{DashPunctuation} & Pd: Punctuation, Dash \\
\haddockdecltt{|} & \haddockdecltt{OpenPunctuation} & Ps: Punctuation, Open \\
\haddockdecltt{|} & \haddockdecltt{ClosePunctuation} & Pe: Punctuation, Close \\
\haddockdecltt{|} & \haddockdecltt{InitialQuote} & Pi: Punctuation, Initial quote \\
\haddockdecltt{|} & \haddockdecltt{FinalQuote} & Pf: Punctuation, Final quote \\
\haddockdecltt{|} & \haddockdecltt{OtherPunctuation} & Po: Punctuation, Other \\
\haddockdecltt{|} & \haddockdecltt{MathSymbol} & Sm: Symbol, Math \\
\haddockdecltt{|} & \haddockdecltt{CurrencySymbol} & Sc: Symbol, Currency \\
\haddockdecltt{|} & \haddockdecltt{ModifierSymbol} & Sk: Symbol, Modifier \\
\haddockdecltt{|} & \haddockdecltt{OtherSymbol} & So: Symbol, Other \\
\haddockdecltt{|} & \haddockdecltt{Space} & Zs: Separator, Space \\
\haddockdecltt{|} & \haddockdecltt{LineSeparator} & Zl: Separator, Line \\
\haddockdecltt{|} & \haddockdecltt{ParagraphSeparator} & Zp: Separator, Paragraph \\
\haddockdecltt{|} & \haddockdecltt{Control} & Cc: Other, Control \\
\haddockdecltt{|} & \haddockdecltt{Format} & Cf: Other, Format \\
\haddockdecltt{|} & \haddockdecltt{Surrogate} & Cs: Other, Surrogate \\
\haddockdecltt{|} & \haddockdecltt{PrivateUse} & Co: Other, Private Use \\
\haddockdecltt{|} & \haddockdecltt{NotAssigned} & Cn: Other, Not Assigned \\
\end{tabulary}\par}
\end{haddockdesc}
\begin{haddockdesc}
\item[\begin{tabular}{@{}l}
instance Bounded GeneralCategory\\instance Enum GeneralCategory\\instance Ix GeneralCategory\\instance Read GeneralCategory\\instance Show GeneralCategory\\instance Eq GeneralCategory\\instance Ord GeneralCategory
\end{tabular}]
\end{haddockdesc}
\begin{haddockdesc}
\item[\begin{tabular}{@{}l}
generalCategory :: Char -> GeneralCategory
\end{tabular}]
{\haddockbegindoc
The Unicode general category of the character. This relies on the
 \haddockid{Enum} instance of \haddockid{GeneralCategory}, which must remain in the
 same order as the categories are presented in the Unicode
 standard.\par
\subsubsection*{\textbf{Examples}}
Basic usage:\par
\begin{quote}
{\haddockverb\begin{verbatim}
>>> generalCategory 'a'
LowercaseLetter

>>> generalCategory 'A'
UppercaseLetter

>>> generalCategory '0'
DecimalNumber

>>> generalCategory '%'
OtherPunctuation

>>> generalCategory '♥'
OtherSymbol

>>> generalCategory '\31'
Control

>>> generalCategory ' '
Space

\end{verbatim}}
\end{quote}}
\end{haddockdesc}
\section{Case conversion}
\begin{haddockdesc}
\item[\begin{tabular}{@{}l}
toUpper :: Char -> Char
\end{tabular}]
{\haddockbegindoc
Convert a letter to the corresponding upper-case letter, if any.
 Any other character is returned unchanged.\par}
\end{haddockdesc}
\begin{haddockdesc}
\item[\begin{tabular}{@{}l}
toLower :: Char -> Char
\end{tabular}]
{\haddockbegindoc
Convert a letter to the corresponding lower-case letter, if any.
 Any other character is returned unchanged.\par}
\end{haddockdesc}
\begin{haddockdesc}
\item[\begin{tabular}{@{}l}
toTitle :: Char -> Char
\end{tabular}]
{\haddockbegindoc
Convert a letter to the corresponding title-case or upper-case
 letter, if any.  (Title case differs from upper case only for a small
 number of ligature letters.)
 Any other character is returned unchanged.\par}
\end{haddockdesc}
\section{Single digit characters}
\begin{haddockdesc}
\item[\begin{tabular}{@{}l}
digitToInt :: Char -> Int
\end{tabular}]
{\haddockbegindoc
Convert a single digit \haddockid{Char} to the corresponding \haddockid{Int}.  This
 function fails unless its argument satisfies \haddockid{isHexDigit}, but
 recognises both upper- and lower-case hexadecimal digits (that
 is, \haddocktt{'0'}..\haddocktt{'9'}, \haddocktt{'a'}..\haddocktt{'f'}, \haddocktt{'A'}..\haddocktt{'F'}).\par
\subsubsection*{\textbf{Examples}}
Characters \haddocktt{'0'} through \haddocktt{'9'} are converted properly to
 \haddocktt{0..9}:\par
\begin{quote}
{\haddockverb\begin{verbatim}
>>> map digitToInt ['0'..'9']
[0,1,2,3,4,5,6,7,8,9]

\end{verbatim}}
\end{quote}
Both upper- and lower-case \haddocktt{'A'} through \haddocktt{'F'} are converted
 as well, to \haddocktt{10..15}.\par
\begin{quote}
{\haddockverb\begin{verbatim}
>>> map digitToInt ['a'..'f']
[10,11,12,13,14,15]

>>> map digitToInt ['A'..'F']
[10,11,12,13,14,15]

\end{verbatim}}
\end{quote}
Anything else throws an exception:\par
\begin{quote}
{\haddockverb\begin{verbatim}
>>> digitToInt 'G'
*** Exception: Char.digitToInt: not a digit 'G'

>>> digitToInt '♥'
*** Exception: Char.digitToInt: not a digit '\9829'

\end{verbatim}}
\end{quote}}
\end{haddockdesc}
\begin{haddockdesc}
\item[\begin{tabular}{@{}l}
intToDigit :: Int -> Char
\end{tabular}]
{\haddockbegindoc
Convert an \haddockid{Int} in the range \haddocktt{0}..\haddocktt{15} to the corresponding single
 digit \haddockid{Char}.  This function fails on other inputs, and generates
 lower-case hexadecimal digits.\par}
\end{haddockdesc}
\section{Numeric representations}
\begin{haddockdesc}
\item[\begin{tabular}{@{}l}
ord :: Char -> Int
\end{tabular}]
{\haddockbegindoc
The \haddockid{fromEnum} method restricted to the type \haddockid{Char}.\par}
\end{haddockdesc}
\begin{haddockdesc}
\item[\begin{tabular}{@{}l}
chr :: Int -> Char
\end{tabular}]
{\haddockbegindoc
The \haddockid{toEnum} method restricted to the type \haddockid{Char}.\par}
\end{haddockdesc}
\section{String representations}
\begin{haddockdesc}
\item[\begin{tabular}{@{}l}
showLitChar :: Char -> ShowS
\end{tabular}]
{\haddockbegindoc
Convert a character to a string using only printable characters,
 using Haskell source-language escape conventions.  For example:\par
\begin{quote}
{\haddockverb\begin{verbatim}
showLitChar '\n' s  =  "\\n" ++ s\end{verbatim}}
\end{quote}}
\end{haddockdesc}
\begin{haddockdesc}
\item[\begin{tabular}{@{}l}
lexLitChar :: ReadS String
\end{tabular}]
{\haddockbegindoc
Read a string representation of a character, using Haskell
 source-language escape conventions.  For example:\par
\begin{quote}
{\haddockverb\begin{verbatim}
lexLitChar  "\\nHello"  =  [("\\n", "Hello")]\end{verbatim}}
\end{quote}}
\end{haddockdesc}
\begin{haddockdesc}
\item[\begin{tabular}{@{}l}
readLitChar :: ReadS Char
\end{tabular}]
{\haddockbegindoc
Read a string representation of a character, using Haskell
 source-language escape conventions, and convert it to the character
 that it encodes.  For example:\par
\begin{quote}
{\haddockverb\begin{verbatim}
readLitChar "\\nHello"  =  [('\n', "Hello")]\end{verbatim}}
\end{quote}}
\end{haddockdesc}